\chapter{标准库速查}
\label{ch:appB}

查标准库时，从要处理的数据和希望得到的结果出发：数值计算、文本转换、容器遍历与外部资源需要不同接口。模块名帮助定位能力，签名说明输入和返回值，失败约定说明调用前提与恢复路径。

本附录用一个组合例程展示查询与调用方法，再按任务提供查阅入口。简短示例不能覆盖每个 API 的全部边界；实际使用还需检查具体类型、是否修改来源、资源归属和平台条件。

\section{先看一个例子}

\par
\medskip
\begin{lstlisting}[style=chali]
import std::math;
import std::text;
import std::iter;

void main()
{
    println("sqrt={}", std::math::sqrt(144.0));
    println("upper={}", std::text::Text::toUpper("chali"));
    int64 sum = 0;
    for (int64 v in std::iter::range(1, 5))
    {
        sum += v;
    }
    println("sum={}", sum);
}
\end{lstlisting}
\par
\medskip

\par
\medskip
\begin{lstlisting}[style=out]
sqrt=12
upper=CHALI
sum=10
\end{lstlisting}
\par
\medskip

\subsection*{代码详解}

三个 import 对应三项不同任务：math 提供平方根，text 提供文本处理，iter 提供区间遍历。限定名让读者知道每个操作来自哪个模块。sqrt(144.0) 得到 12，toUpper 把 chali 转成 CHALI，分别填入两条输出。

最后 range(1,5) 产生 1、2、3、4，循环外的 sum 持续累加得到 10。数学与文本结果直接打印，区间结果则先经过循环归约；import 只是使模块可用，不会自动执行这些业务操作。查库时可按输入类型、返回值和失败方式定位接口，再组合成自己的流程。

三行分别来自三个模块：数学、文本、迭代。注意 \texttt{range(\allowbreak{}1,\allowbreak{} 5)\allowbreak{}} 是左闭右开 （拿到 1、2、3、4），所以 \texttt{sum} 是 10。其他 API 的端点、索引和失败方式按各自参考页核对。

\section{核心模块 \texttt{std:\allowbreak{}:\allowbreak{}lang}（不用 import）}

这几样任何程序里都能直接用：

{\small\setlength{\tabcolsep}{4pt}
\begin{longtable}{>{\raggedright\arraybackslash}p{\dimexpr 0.088\linewidth-2\tabcolsep\relax}>{\raggedright\arraybackslash}p{\dimexpr 0.892\linewidth-2\tabcolsep\relax}}
\toprule
能力 & 用法 \\
\midrule
\endhead
输出 & \texttt{println(\allowbreak{}fmt,\allowbreak{} …)\allowbreak{}}、\texttt{format(\allowbreak{}fmt,\allowbreak{} …)\allowbreak{}} \\
断言 & \texttt{assert(\allowbreak{}cond)\allowbreak{}} \\
取长度 & \texttt{len(\allowbreak{}x)\allowbreak{}}（字符串 / 容器） \\
字符串 & \texttt{std:\allowbreak{}:\allowbreak{}lang:\allowbreak{}:\allowbreak{}String} \\
容器 & \texttt{List<\allowbreak{}T>\allowbreak{}}、\texttt{Map<\allowbreak{}K,\allowbreak{} V>\allowbreak{}}、\texttt{Set<\allowbreak{}T>\allowbreak{}}、\texttt{Pair<\allowbreak{}K,\allowbreak{} V>\allowbreak{}} \\
可选值与错误 & \texttt{Option<\allowbreak{}T>\allowbreak{}}、\texttt{Result<\allowbreak{}T,\allowbreak{} E>\allowbreak{}}；标准库错误实现 \texttt{Error.\allowbreak{}to\allowbreak{}String(\allowbreak{})\allowbreak{}}，文本错误使用 \texttt{Message\allowbreak{}Error}，配合模式匹配、查询和取值方法 \\
类型擦除盒子 & \texttt{Any}（第 13 章的运行期多态） \\
数学常数 & \texttt{Math:\allowbreak{}:\allowbreak{}PI}、\texttt{Math:\allowbreak{}:\allowbreak{}E} \\
大数 & \texttt{Big\allowbreak{}Integer}、\texttt{Big\allowbreak{}Decimal} \\
输入 & \texttt{input\allowbreak{}Int(\allowbreak{})\allowbreak{}}、\texttt{input\allowbreak{}UInt(\allowbreak{})\allowbreak{}}、\texttt{input\allowbreak{}Float(\allowbreak{})\allowbreak{}}、\texttt{input\allowbreak{}Bool(\allowbreak{})\allowbreak{}}、\texttt{input\allowbreak{}Char(\allowbreak{})\allowbreak{}}、\texttt{input\allowbreak{}Big\allowbreak{}Integer(\allowbreak{})\allowbreak{}}、\texttt{input\allowbreak{}Big\allowbreak{}Decimal(\allowbreak{})\allowbreak{}} \\
整行输入 & \texttt{read\allowbreak{}Line(\allowbreak{})\allowbreak{}}，保留空格，不含 LF/CRLF \\
多行输入 & \texttt{read\allowbreak{}Lines(\allowbreak{}int32 N)\allowbreak{}}，LF 连接，EOF 提前结束 \\
\bottomrule
\end{longtable}
}

完整类型、接口与内建入口见\href{https://github.com/rabbitc999/ChaLi/blob/master/doc/library/std-lang.md}{核心库参考}。其余模块按需显式 \texttt{import}。

\section{按能力域速查}

\subsection*{数学与数值}

{\small\setlength{\tabcolsep}{4pt}
\begin{longtable}{>{\raggedright\arraybackslash}p{\dimexpr 0.262\linewidth-2\tabcolsep\relax}>{\raggedright\arraybackslash}p{\dimexpr 0.718\linewidth-2\tabcolsep\relax}}
\toprule
模块 & 能干什么 \\
\midrule
\endhead
\href{https://github.com/rabbitc999/ChaLi/blob/master/doc/library/std-math.md}{\texttt{std:\allowbreak{}:\allowbreak{}math}} & \texttt{sin} \texttt{cos} \texttt{sqrt} 等常用函数；\texttt{abs} \texttt{min} \texttt{max} 是泛型入口，\texttt{min}/\texttt{max} 可以给任意多个实参（常数见 \texttt{std:\allowbreak{}:\allowbreak{}lang:\allowbreak{}:\allowbreak{}Math}） \\
\href{https://github.com/rabbitc999/ChaLi/blob/master/doc/library/std-numerics.md}{\texttt{std:\allowbreak{}:\allowbreak{}numerics}} & 标量数值方法：导数、辛普森积分、二分、RK4、牛顿迭代 \\
\href{https://github.com/rabbitc999/ChaLi/blob/master/doc/library/std-numtheory.md}{\texttt{std:\allowbreak{}:\allowbreak{}numtheory}} & 整数数论：模块级泛型 \texttt{gcd} \texttt{lcm}（任意多个实参）、\texttt{add\allowbreak{}Mod} \texttt{mul\allowbreak{}Mod} \texttt{pow\allowbreak{}Mod} \texttt{mod\allowbreak{}Inverse}，以及确定性的 \texttt{is\allowbreak{}Prime}（Miller-Rabin） \\
\href{https://github.com/rabbitc999/ChaLi/blob/master/doc/library/std-complex.md}{\texttt{std:\allowbreak{}:\allowbreak{}complex}} & 复数 \texttt{Complex<\allowbreak{}T>\allowbreak{}} 及对应算术运算 \\
\href{https://github.com/rabbitc999/ChaLi/blob/master/doc/library/std-linalg.md}{\texttt{std:\allowbreak{}:\allowbreak{}linalg}} & 稠密向量与矩阵：\texttt{lu\allowbreak{}Decompose} \texttt{symmetric\allowbreak{}Eigen} \texttt{svd} \texttt{jordan2} \texttt{diagonalize2} \\
\href{https://github.com/rabbitc999/ChaLi/blob/master/doc/library/std-random.md}{\texttt{std:\allowbreak{}:\allowbreak{}random}} & 确定性随机数（LCG）与熵播种；句柄型状态，实现 \texttt{Drop} \\
\href{https://github.com/rabbitc999/ChaLi/blob/master/doc/library/std-signal.md}{\texttt{std:\allowbreak{}:\allowbreak{}signal}} & DFT / FFT 与逆变换；基 2 的 FFT 只接受 2 的幂长度 \\
\href{https://github.com/rabbitc999/ChaLi/blob/master/doc/library/std-cas.md}{\texttt{std:\allowbreak{}:\allowbreak{}cas}} & \texttt{Expr} 精确符号表达式：解析、节点遍历、化简、求导、积分、替换、极限和 Taylor 展开 \\
\bottomrule
\end{longtable}
}

\subsection*{文本与数据格式}

{\small\setlength{\tabcolsep}{4pt}
\begin{longtable}{>{\raggedright\arraybackslash}p{\dimexpr 0.329\linewidth-2\tabcolsep\relax}>{\raggedright\arraybackslash}p{\dimexpr 0.651\linewidth-2\tabcolsep\relax}}
\toprule
模块 & 能干什么 \\
\midrule
\endhead
\href{https://github.com/rabbitc999/ChaLi/blob/master/doc/library/std-text.md}{\texttt{std:\allowbreak{}:\allowbreak{}text}} & \texttt{Text:\allowbreak{}:\allowbreak{}find}、\texttt{split}、\texttt{join}、\texttt{replace}、\texttt{trim}、\texttt{to\allowbreak{}Upper}、\texttt{pad\allowbreak{}Start} 及字节级 KMP \\
\href{https://github.com/rabbitc999/ChaLi/blob/master/doc/library/std-utf8.md}{\texttt{std:\allowbreak{}:\allowbreak{}utf8}} & UTF-8 有效性检查、码点解码与编码、按码点计数和截取 \\
\href{https://github.com/rabbitc999/ChaLi/blob/master/doc/library/std-encoding.md}{\texttt{std:\allowbreak{}:\allowbreak{}encoding}} & 十六进制、Base64 和 URL 百分号编码与解码 \\
\href{https://github.com/rabbitc999/ChaLi/blob/master/doc/library/std-csv.md}{\texttt{std:\allowbreak{}:\allowbreak{}csv}} & CSV 解析与写出，处理字段分隔、引号和换行 \\
\href{https://github.com/rabbitc999/ChaLi/blob/master/doc/library/std-json.md}{\texttt{std:\allowbreak{}:\allowbreak{}json}} & \texttt{parse} \ensuremath{\rightarrow} \texttt{Result<\allowbreak{}Json\allowbreak{}Value,\allowbreak{} Error>\allowbreak{}}；再 \texttt{is\allowbreak{}Int} \texttt{is\allowbreak{}Object} \texttt{len} \texttt{as\allowbreak{}Int} \texttt{as\allowbreak{}String} 取值 \\
\href{https://github.com/rabbitc999/ChaLi/blob/master/doc/library/std-toml.md}{\texttt{std:\allowbreak{}:\allowbreak{}toml}} & \texttt{Toml\allowbreak{}Value} 文档访问、解析和保留格式的转储，以及反射编解码 \\
\href{https://github.com/rabbitc999/ChaLi/blob/master/doc/library/std-regex.md}{\texttt{std:\allowbreak{}:\allowbreak{}regex}} & \texttt{Regex:\allowbreak{}:\allowbreak{}compile}，整串匹配、包含、查找、捕获和替换；编译失败返回 \texttt{Regex\allowbreak{}Error} \\
\href{https://github.com/rabbitc999/ChaLi/blob/master/doc/library/std-parse.md}{\texttt{std:\allowbreak{}:\allowbreak{}parse}} & 字符串 \ensuremath{\rightarrow} 整数，失败给 \texttt{Parse\allowbreak{}Error} \\
\href{https://github.com/rabbitc999/ChaLi/blob/master/doc/library/std-binary.md}{\texttt{std:\allowbreak{}:\allowbreak{}binary}} & 定宽解码、字节序处理与边界检查 \\
\href{https://github.com/rabbitc999/ChaLi/blob/master/doc/library/std-path.md}{\texttt{std:\allowbreak{}:\allowbreak{}path}} & 路径拼接、取文件名、取扩展名、规范化 \\
\href{https://github.com/rabbitc999/ChaLi/blob/master/doc/library/std-compress.md}{\texttt{std:\allowbreak{}:\allowbreak{}compress}} & 文本压缩 / 解压；压缩结果是二进制 \texttt{List<\allowbreak{}uint8>\allowbreak{}} \\
\bottomrule
\end{longtable}
}

\subsection*{数据结构、通用算法与图像}

{\small\setlength{\tabcolsep}{4pt}
\begin{longtable}{>{\raggedright\arraybackslash}p{\dimexpr 0.394\linewidth-2\tabcolsep\relax}>{\raggedright\arraybackslash}p{\dimexpr 0.586\linewidth-2\tabcolsep\relax}}
\toprule
模块 & 能干什么 \\
\midrule
\endhead
\href{https://github.com/rabbitc999/ChaLi/blob/master/doc/library/std-collections.md}{\texttt{std:\allowbreak{}:\allowbreak{}collections}} & \texttt{Deque}、\texttt{Queue}、当前上下文的 \texttt{Channel}、\texttt{Fixed\allowbreak{}Vector} 和 \texttt{Ring\allowbreak{}Buffer} \\
\href{https://github.com/rabbitc999/ChaLi/blob/master/doc/library/std-linkedlist.md}{\texttt{std:\allowbreak{}:\allowbreak{}linkedlist}} & 双向链表，头尾操作、插入、删除、查找和遍历 \\
\href{https://github.com/rabbitc999/ChaLi/blob/master/doc/library/std-tree.md}{\texttt{std:\allowbreak{}:\allowbreak{}tree}} & 二叉搜索树、红黑树、最小堆和普通树；删除、遍历、校验与快照 \\
\href{https://github.com/rabbitc999/ChaLi/blob/master/doc/library/std-graph.md}{\texttt{std:\allowbreak{}:\allowbreak{}graph}} & \texttt{Graph}、\texttt{Union\allowbreak{}Find}，BFS、DFS、拓扑排序、最短路、环检测和最小生成树 \\
\href{https://github.com/rabbitc999/ChaLi/blob/master/doc/library/std-algorithms.md}{\texttt{std:\allowbreak{}:\allowbreak{}algorithms}} & 排序、查找、折叠、映射、过滤、计数和聚合；回调遵循函数值的借用与生命周期规则 \\
\href{https://github.com/rabbitc999/ChaLi/blob/master/doc/library/std-image.md}{\texttt{std:\allowbreak{}:\allowbreak{}image}} & \texttt{Image}、RGBA 像素，BMP/PGM/PPM 读写、颜色转换、卷积、池化和大津法二值化 \\
\bottomrule
\end{longtable}
}

\texttt{std:\allowbreak{}:\allowbreak{}collections:\allowbreak{}:\allowbreak{}Channel} 是当前执行上下文的值队列。跨线程传递使用\texttt{std:\allowbreak{}:\allowbreak{}concurrency:\allowbreak{}:\allowbreak{}Channel<\allowbreak{}T>\allowbreak{}} 或 \texttt{std:\allowbreak{}:\allowbreak{}message}，两类队列的共享语义不同。

\subsection*{反射与模型编解码}

{\small\setlength{\tabcolsep}{4pt}
\begin{longtable}{>{\raggedright\arraybackslash}p{\dimexpr 0.367\linewidth-2\tabcolsep\relax}>{\raggedright\arraybackslash}p{\dimexpr 0.613\linewidth-2\tabcolsep\relax}}
\toprule
模块 & 能干什么 \\
\midrule
\endhead
\href{https://github.com/rabbitc999/ChaLi/blob/master/doc/library/std-reflect.md}{\texttt{std:\allowbreak{}:\allowbreak{}reflect}} & 查询类型、字段、方法、注解和稳定类型 ID；受反射可见性约束的读写与动态调用 \\
\href{https://github.com/rabbitc999/ChaLi/blob/master/doc/library/std-serialization.md}{\texttt{std:\allowbreak{}:\allowbreak{}serialization}} & \texttt{serialize<\allowbreak{}T>\allowbreak{}}、\texttt{deserialize<\allowbreak{}T>\allowbreak{}}，字节流保存版本、稳定类型 ID 和长度，恢复时严格检查模型类型 \\
\bottomrule
\end{longtable}
}

JSON/TOML 的模型编解码见对应格式模块，数据库模型与动态 SQL 见 \texttt{std:\allowbreak{}:\allowbreak{}sql}。源码编译默认启用反射；关闭本次元数据输出时显式使用 \texttt{-\allowbreak{}-\allowbreak{}no-\allowbreak{}reflection}。执行使用 checked AOT 或导出原生代码。

\subsection*{并发与协作}

{\small\setlength{\tabcolsep}{4pt}
\begin{longtable}{>{\raggedright\arraybackslash}p{\dimexpr 0.365\linewidth-2\tabcolsep\relax}>{\raggedright\arraybackslash}p{\dimexpr 0.615\linewidth-2\tabcolsep\relax}}
\toprule
模块 & 能干什么 \\
\midrule
\endhead
\href{https://github.com/rabbitc999/ChaLi/blob/master/doc/library/std-thread.md}{\texttt{std:\allowbreak{}:\allowbreak{}thread}} & 隔离线程 \texttt{Thread}、取消、结果、线程组和固定线程池 \\
\href{https://github.com/rabbitc999/ChaLi/blob/master/doc/library/std-message.md}{\texttt{std:\allowbreak{}:\allowbreak{}message}} & 跨线程的消息 \\
\href{https://github.com/rabbitc999/ChaLi/blob/master/doc/library/std-concurrency.md}{\texttt{std:\allowbreak{}:\allowbreak{}concurrency}} & 有界跨线程 \texttt{Channel<\allowbreak{}T>\allowbreak{}}、显式共享槽 \texttt{Shared<\allowbreak{}T>\allowbreak{}} 和 \texttt{Atomic\allowbreak{}Int64}（普通全局变量仍隔离） \\
\href{https://github.com/rabbitc999/ChaLi/blob/master/doc/library/std-event.md}{\texttt{std:\allowbreak{}:\allowbreak{}event}} & 显式事件与定时源；通过 receive 接收，周期积压合并 \\
\bottomrule
\end{longtable}
}

\begin{tip}{提示}
线程执行、通道通信和显式事件分别由 std::thread、std::concurrency 和 std::event 提供。
\end{tip}

\subsection*{系统与环境}

{\small\setlength{\tabcolsep}{4pt}
\begin{longtable}{>{\raggedright\arraybackslash}p{\dimexpr 0.263\linewidth-2\tabcolsep\relax}>{\raggedright\arraybackslash}p{\dimexpr 0.717\linewidth-2\tabcolsep\relax}}
\toprule
模块 & 能干什么 \\
\midrule
\endhead
\href{https://github.com/rabbitc999/ChaLi/blob/master/doc/library/std-io.md}{\texttt{std:\allowbreak{}:\allowbreak{}io}} & 文件读写；错误类型 \texttt{Io\allowbreak{}Error}（\texttt{Open\allowbreak{}Failed} \texttt{Read\allowbreak{}Failed} \texttt{Write\allowbreak{}Failed} \texttt{Seek\allowbreak{}Failed} \texttt{Closed}） \\
\href{https://github.com/rabbitc999/ChaLi/blob/master/doc/library/std-net.md}{\texttt{std:\allowbreak{}:\allowbreak{}net}} & 独占的 \texttt{Tcp\allowbreak{}Socket}、\texttt{Tcp\allowbreak{}Listener}，连接、监听、收发和 HTTP 客户端；HTTP 客户端不提供 TLS \\
\href{https://github.com/rabbitc999/ChaLi/blob/master/doc/library/std-sql.md}{\texttt{std:\allowbreak{}:\allowbreak{}sql}} & \texttt{Database}、事务与模型操作；连接 MySQL / SQLite / PostgreSQL / SQL Server。参数绑定与后端能力分别核对，MySQL 文本使用 quote \\
\href{https://github.com/rabbitc999/ChaLi/blob/master/doc/library/std-process.md}{\texttt{std:\allowbreak{}:\allowbreak{}process}} & \texttt{Process}，宿主进程创建、命令管道、等待与结果读取 \\
\href{https://github.com/rabbitc999/ChaLi/blob/master/doc/library/std-clock.md}{\texttt{std:\allowbreak{}:\allowbreak{}clock}} & 单调时钟 \texttt{steady\_\allowbreak{}clock}，只用来做差值，别当挂钟用 \\
\href{https://github.com/rabbitc999/ChaLi/blob/master/doc/library/std-datetime.md}{\texttt{std:\allowbreak{}:\allowbreak{}datetime}} & \texttt{Date\allowbreak{}Time}、UTC 公历日期、Unix 秒、格式化和解析；不表示闰秒 \\
\href{https://github.com/rabbitc999/ChaLi/blob/master/doc/library/std-cli.md}{\texttt{std:\allowbreak{}:\allowbreak{}cli}} & 长短选项、组合短选项、位置参数、\texttt{-\allowbreak{}-\allowbreak{}} 分隔、参数查询和帮助生成 \\
\href{https://github.com/rabbitc999/ChaLi/blob/master/doc/library/std-memory.md}{\texttt{std:\allowbreak{}:\allowbreak{}memory}} & \texttt{Span<\allowbreak{}T>\allowbreak{}} 借用视图与 \texttt{Memory\allowbreak{}Stats}；缓冲区用 \texttt{new} 分配，视图受底层存储有效期约束 \\
\bottomrule
\end{longtable}
}

\subsection*{迭代}

{\small\setlength{\tabcolsep}{4pt}
\begin{longtable}{>{\raggedright\arraybackslash}p{\dimexpr 0.342\linewidth-2\tabcolsep\relax}>{\raggedright\arraybackslash}p{\dimexpr 0.638\linewidth-2\tabcolsep\relax}}
\toprule
模块 & 能干什么 \\
\midrule
\endhead
\href{https://github.com/rabbitc999/ChaLi/blob/master/doc/library/std-iter.md}{\texttt{std:\allowbreak{}:\allowbreak{}iter}} & \texttt{range}、\texttt{map} / \texttt{filter} / \texttt{collect} 惰性链，以及 \texttt{Iter\allowbreak{}Next<\allowbreak{}T>\allowbreak{}} 接口 \\
\bottomrule
\end{longtable}
}

\href{https://github.com/rabbitc999/ChaLi/blob/master/doc/library/std-probe.md}{\texttt{std:\allowbreak{}:\allowbreak{}probe}}用于最小宿主集成探针与验证，不作为应用功能模块。

\section{三条容易踩的约定}

一、资源按类型的生命周期管理。 拥有型对象按 \texttt{Drop} 约定自动释放，提前释放使用类型提供的接口。不能假定所有对象都有相同的 \texttt{close(\allowbreak{})\allowbreak{}} 或重复关闭行为。独占文件、连接和进程不能随意复制；转移所有权写 \texttt{move value}。复制是否深复制、借用何时失效，查对应类型的参考页。

二、压缩结果是二进制，不是文本。 \texttt{std:\allowbreak{}:\allowbreak{}compress} 的输出含 \texttt{\textbackslash{}0} 字节，类型是\texttt{List<\allowbreak{}uint8>\allowbreak{}}；别拿它当字符串打印，解压出来的才是 \texttt{std:\allowbreak{}:\allowbreak{}lang:\allowbreak{}:\allowbreak{}String}。

三、字节与码点分清。 \texttt{String} 的长度和切片端点按字节计算，切片可以截取多字节编码的一部分；普通字符串下标仍要求 UTF-8 字符边界。按 Unicode 码点处理时使用显式码点接口或 \texttt{std:\allowbreak{}:\allowbreak{}utf8}。

实际文件、网络、数据库和线程能力以运行期胶水、系统库及构建配置为准。模块是否可在某个目标使用，还取决于宿主服务、链接资产和外部库；按 \href{https://github.com/rabbitc999/ChaLi/blob/master/doc/38-aot.md}{AOT 参考} 与模块页分别验证。

\section{细节去哪儿查}

{\small\setlength{\tabcolsep}{4pt}
\begin{longtable}{>{\raggedright\arraybackslash}p{\dimexpr 0.318\linewidth-2\tabcolsep\relax}>{\raggedright\arraybackslash}p{\dimexpr 0.662\linewidth-2\tabcolsep\relax}}
\toprule
你想查 & 去哪 \\
\midrule
\endhead
某个函数的确切签名、默认值、返回与错误 & \href{https://github.com/rabbitc999/ChaLi/blob/master/doc/library/README.md}{库函数 API 参考} 与\href{https://github.com/rabbitc999/ChaLi/blob/master/doc/library/symbols.md}{全部函数索引} \\
按能力域通读一遍标准库 & \href{https://github.com/rabbitc999/ChaLi/blob/master/doc/07-stdlib-overview.md}{07 标准库概览} \\
标准库源码（每个模块头部都有能力说明） & \href{https://github.com/rabbitc999/ChaLi/tree/master/source/std}{\texttt{source/\allowbreak{}std/\allowbreak{}}} 下的 \texttt{source/\allowbreak{}std/\allowbreak{}<\allowbreak{}名字>\allowbreak{}.\allowbreak{}chl} \\
宿主侧绑定（\texttt{std:\allowbreak{}:\allowbreak{}*} 的真正实现） & \href{https://github.com/rabbitc999/ChaLi/tree/master/source/glue}{\texttt{source/\allowbreak{}glue/\allowbreak{}}} 下按功能组织的 \texttt{std\_\allowbreak{}*} 模块 \\
可直接运行的标准库样例 & \href{https://github.com/rabbitc999/ChaLi/tree/master/example/samples}{\texttt{example/\allowbreak{}samples/\allowbreak{}}} \\
\bottomrule
\end{longtable}
}

模块声明、宿主胶水和编译器注入实现按职责组织，文件名不必与模块名逐字对应；查询具体实现来源时以 API 索引中的源码位置为准。

\section{组合与验证}

选择模块后检查输入类型、返回的 Result／Option、资源关闭及目标平台条件。先运行最小正常例程，再验证一个失败路径；跨模块组合时重新检查所有权与借用关系。数据库组合示例见\hyperref[ch:25]{第 25 章}，接口错误匹配见\hyperref[ch:16]{第 16 章}。
